\documentclass[../tufte-style.tex]{standalone}

\begin{document}

\date{February 23, 2026}

\begin{frame}{Chain rule}

\begin{block}{Definition of the chain rule}

The chain rule allow us to compute the derivative of a composition of functions \(f(x) = \left( h\circ g \right)(x)=h(g(x))\).
The rule goes as follows,

\vspace{2em}

\begin{minipage}{0.45\textwidth}
    \centering
     Leibniz notation
    \begin{gather*}
        \dv{f(x)}{x} = \left[\dv{h(x)}{x}\circ g(x)\right]\dv{g(x)}{x}
    \end{gather*}
\end{minipage}
\hfill
\begin{minipage}{0.45\textwidth}
    \centering
    Lagrange notation
    \begin{gather*}
        f'(x) = h'(g(x))\cdot g'(x) 
    \end{gather*}
\end{minipage}

\end{block}

\end{frame}

\begin{frame}{Chain rule}
\begin{examples}
Let's compute the derivative of \(f(x) = \ln\qty(\cos\qty(x))\).

\begin{align*}
    f'(x) &= \ln'\qty(\cos\qty(x))\cdot\cos'\qty(x) \\
          &= \frac{1}{\cos\qty(x)}\cdot-\sin\qty(x) \\
          &= -\tan(x)
\end{align*}

\end{examples}
\end{frame}

\begin{frame}{Chain rule}
\begin{examples}
Let's compute the derivative of \(f(x) = \cos\qty(2\pi x-\pi)\).

\begin{align*}
    f'(x) &= \cos'\qty(2\pi x-\pi)\cdot\qty(2\pi x-\pi)' \\
          &= -\sin\qty(2\pi x-\pi)\cdot\qty(2\pi) \\
          &= -2\pi\sin\qty(2\pi x-\pi)
\end{align*}

\end{examples}
\end{frame}

\begin{frame}{Chain rule}
\begin{examples}
Let's compute the derivative of \(f(x) = e^{-\frac{x^2}{4}}\).

\begin{align*}
    f'(x) &= \left[\dv{e^{x}}{x}\circ-\frac{x^2}{4}\right]\cdot\qty(-\frac{x^2}{4})' \\
          &= e^{-\frac{x^2}{4}}\cdot\qty(-\frac{x}{2}) \\
          &= -\frac{x}{2}e^{-\frac{x^2}{4}}
\end{align*}

\end{examples}
\end{frame}

\begin{frame}{Chain rule}

\begin{block}{Class Activity}
    Compute the derivative of the following functions using the chain rule:
    \begin{gather*}
        f(x) = \ln(e^x),\qquad g(x) = e^{\cos(x)}
    \end{gather*}
\end{block}

\end{frame}



\end{document}
